\section{All counts, finite implementation, and the guarantee}\label{simp:finite}

If no positive singleton perfect matching exists and $|M|\ge n$, reserve
any $n$ items and give one to each agent; then $\OPT=0$ and the guarantee
is immediate. If $|M|<n$, any allocation suffices. Below assume a positive
reservation, so every logarithmic comparison is finite.

\subsection{No further head counts or scalar cases}

Fix one row and its residual comparator. For $h=|H|\ge4$, nonnegative
monotone submodularity gives
\[
 \sum_{I\subseteq H,\ |I|=4}f(O\cup I)
                         \ge {h-1\choose3}f(O\cup H).
\]
Expand into ordered nonnegative marginal increments: each item occurs in
exactly ${h-1\choose3}$ subsets and its marginal in a smaller set is at
least its full ordered marginal. Extra copies of $f(O)\ge0$ help.
Since $\widetilde P_j\ge0$, exponentiating and summing the proved four-head
price inequalities yields
$\mathcal E(H)+\Psi(H)\le\log h$. A row-independent $+b$ paid by every
four-subset is retained once in this inequality.
For B this suffices since $\log h\le\log4+(h-4)\lambda_\star$.
For A the remaining arithmetic charge is at most $(6h/25)\zeta$.
For every real $h\ge5$,
\[
 (h-4)\lambda_\star-\log(h/4)-(6h/25)\zeta>0:
\]
its value at five exceeds $.00782988$ and its derivative is at least
$\lambda_\star-1/5-(6/25)\zeta>.03171853$.
Thus \eqref{eq:integration-scalar-interface} holds for all head counts.

Fix an optimal allocation and assign it any unused reserved items.
Its reserved subsets $H_i$ partition $H_0$, so $\sum_i|H_i|=n$.
Let $J=\sum_i\log u_i+\sum_j\pi_j$ and let $S$ be the logarithmic
score after the single sampled repair at an iterate. The graph theorem gives
$\mathbb E[S\mid x]\ge J-\sum_i\ell_i(H_i)$.
Summing \eqref{eq:integration-scalar-interface} cancels both
$\lambda_\star$ and the HOME transfer: each reserved item appears once
as a comparator head and once as a home column. Hence
\begin{equation}\label{simp:iterate-bound}
 n\log\OPT\le\mathbb E[S\mid x]+nU_\star+\sum_i\mathcal B_i.
\end{equation}
The comparator is used only in this proof. All prices, means and credits
come from the same auxiliary solve.

On active normalized rows $u\ge\Gamma\alpha\ge\theta_0/d$ and
$0\le\mu\le1$. Thus rewards are uniformly bounded by $d/\theta_0$.
For $N\le nd$ active coordinates, a bound $G_0$ on implemented rewards,
and $T$ projected steps of size $1/(G_0\sqrt T)$, the squared-distance
identity gives, for every fixed residual comparator,
\[
 \frac1T\sum_{q=1}^T\sum_i\mathcal B_i^{(q)}\le\frac{NG_0}{\sqrt T}
\]
in the exact model. This uses only the radial inequality, not continuity
of the root or the selected LP solution. Select the iterate uniformly
and perform the single repair and common rounding there.
\eqref{simp:iterate-bound} then gives
$\mathbb E\log\NSW\ge\log\OPT-U_\star-\text{regret}/n$,
before finite errors. No repair at discarded iterates is necessary.

\subsection{A uniform finite sensitivity bound}

All sampled residual statistics have range at most one after subtracting
known physical baselines. Let
\[
 \kappa_0=\min\{1,c_0/(4b_s)\},\quad L_0=\lceil2/\kappa_0\rceil,
 \quad r_0=(\kappa_0/4)^{L_0},\quad
 K_f=\left\lceil\frac{10^6(1+r_0^{-1})^2}{\theta_0^4}\right\rceil.
\]
This fixed integer bounds the needed score sensitivities by $K_fd$.
The private log score has sensitivity $1/(\theta_0\alpha)$.
Affine and HOME numerators have bounded mean derivatives. The only
nontrivial case is the joint branch with floor $m-b_s\alpha$.
Its activation implies $F\ge\alpha$. At the evaluation point, a retained
finite physical entry has true value at least $e^{-\eta}c_0F\ge c_0F/2$.
Follow it only until a further halving. With
$r=(m-b_s\alpha)/(b_s\alpha)$, its normalized value satisfies
\[
 T/(b_s\alpha)=\min_{0\le q\le1}
       [1-q+r^{1-q}(1+r)^q]
       \le\kappa_0/2+2r^{\kappa_0/2}\quad(0<r\le1).
\]
Thus $r\ge(\kappa_0/4)^{2/\kappa_0}\ge r_0$ while the normalized
value is at least $\kappa_0$. Differentiation of the variational expression
bounds its mean derivative by $1+1/r_0$. The other joint branch has
mean derivative one and remains on that branch when its mean decreases.
All stored positive scores start above $\theta_0\alpha$; before halving,
their log sensitivity is at most
$2(1+r_0^{-1})/(\theta_0\alpha)\le K_fd$.
The errors below are smaller than a possible halving displacement.
Physical entries dominated by the private option are simply discarded;
this argument does not reintroduce physical guard labels.

For confidence means, subtract baselines first and clamp the upper
contraction mean to the upper raw mean. Evaluate monotone score formulas
at upper mean endpoints and enclose their logarithms upward. The rational
LP is therefore feasible for every true ordinary/private constraint.
The coherent NEW weights, lower-endpoint eligibility, final rectangular
classification and shaved repair thresholds are exactly those of
Section~\ref{simp:algorithm-interface}. They retain every truly eligible
NEW constraint, every positive supported rational mass and the broad
minorant condition. Supported score overstatement costs at most $2\eta$
per row, the false-A strip costs less than $9\eta$, and the upper-root
error costs at most $(32\beta/c_0)\sqrt\eta$. The entire new loss is
less than $64\sqrt\eta$ per row, including critical equality points.

For clarity, the root enclosure is computable without a stability
assumption. Reparameterize $Q=st$, $s\in[0,e^{-1}]$; the credit objective
is $6$-Lipschitz in $s$ on $t\le3$ and changes by at most
$2|\Delta t|+|\Delta D|$ with its normalized means. A mesh of width
$\eta/100$ and enclosing elementary evaluations give a credit/target
error at most $16\eta$. The private target lies in
\[
 \tfrac13\min(v-w_0,P_c)\le b\le
 \tfrac13\min(v-w_0+\eta,P_c).
\]
To use the uncertain admissible endpoint, first enclose it by
$\bar\mu_-:=\min(1,1/t^+)\le\bar\mu\le
\bar\mu_+:=\min(1,1/t^-)$, with the value one at a zero denominator.
These endpoints differ by at most $\eta$. If the lower enclosure of
$H(\bar\mu_-)-b$ is nonpositive, return $\bar\mu_+$ as a root upper
bound; the quadratic separation gives error at most $8\sqrt\eta+\eta$.
Otherwise bisect within the known increasing interval $(0,\bar\mu_-]$.
At an unresolved sign, the true residual has magnitude at most $32\eta$,
so the root lies within $8\sqrt\eta$ of that midpoint. Use the upper
end of that bracket, also bounded by $\bar\mu_-$, or stop at width $\eta$.
Finally clip upward by the upper enclosure of the floor and at one.
The resulting upper multiplier differs from $\mu$ by at most
$32\sqrt\eta$; monotonicity beyond $\bar\mu$ is never assumed.
The normalized radial coefficient is at least $-\beta t$, so upward
root error costs at most the stated uniform amount even if the
computed multiplier slightly exceeds the uncertain $1/t$ endpoint.

\subsection{One explicit budget and every-history feasibility}

If there are no active residual coordinates, return the reservation,
which is optimal; otherwise $d,N\ge1$ in the formulas below.
For a required logarithmic tolerance $\xi>0$, put
$\xi_0=\min(\xi,1)$, $\epsilon_0=\xi_0/10^6$,
$C_f=10^6(K_f+1)^2$, and choose
\begin{equation}\label{simp:precision}
 \eta=\min\left\{10^{-10},
       \left(\frac{\epsilon_0}{10^6C_f(d+1)(N+1)}\right)^4\right\},\qquad
 \delta_m=\frac{\eta}{100C_f(d+1)(N+1)}.
\end{equation}
Use $\delta_m$ for primitive online mean and elementary-function errors,
then enclose scores and normalized means at scale $\eta$.
Division by $u\ge\theta_0/d$ and the preceding $K_fd$ sensitivity are
absorbed by \eqref{simp:precision}. All elementary logarithms are evaluated
to directed rational accuracy; large scores are handled logarithmically.
Take $G_0=4C_fd$ and $T$ the smallest power of four with
$\sqrt T\ge100NG_0/\epsilon_0$. Project with step
$\tau=1/(G_0\sqrt T)$ and reset coordinates downward to mesh at most
$\tau\epsilon_0/(100N)$. Downward resets preserve capacities and add at
most $2Nh+Nh^2$ to squared distance at mesh $h$; the corresponding
regret, reward and reset terms cost at most $4n\epsilon_0$.
For contraction selection omit $p_{ij}>2$ on normalized active rows,
since their bonus is negative, and maximize the remaining estimates
alongside an exact zero dummy. Its true bonus loss is at most
$6\delta_m$ and is absorbed by the same budget.

For the common forest engine use normalized derivative accuracy
$a=\epsilon_0/[10000C_f(d+1)(N+1)]$ and downward mesh at most
$a/(d+1)$. Minorant rows are normalized by dividing their variable
range by two. The explicit cycle/pair losses in Lemma~\ref{simp:forest}
then cost less than $n\epsilon_0$ in log score. A range-one estimate
with error $a'$ and failure allowance $p'$ uses
$\lceil2(a')^{-2}\log(2/p')\rceil$ samples; the standard exponential
moment bound $\mathbb Ee^{s(X-\mathbb EX)}\le e^{s^2/8}$ proves this
also conditionally on every adaptive history.
Dyadic product sampling with $b$ bits changes means by at most $d2^{-b}$.
The radial distribution is generated from a dyadic uniform midpoint $v$
via $s=\log(1+(e-1)v)$, whose derivative is below two.
These biases fit in the primitive estimation allowances.

Retain a private weight at least a downward enclosure of
$\log(\theta_0\alpha)$ with error at most $1/2$ on every history;
clip statistics and rewards to their known ranges. If a selected repair
fails a mandatory eligibility or feasibility check, return the original
reservation. The strict gates exclude this event on accurate histories.
All matchings, flows and forest outcomes are feasible even on inaccurate
histories, and the final matching can always use the positive reservation.
Let $\Lambda_f$ be a rational upper bound for
$\log[\prod_i f_i(M)/\prod_i p_{ii}]/n$, and let
$M_f\ge1+\log(4C_f(d+2)/\theta_0)$ be rational.
The latter bounds the loss of a bad base row: private/ordinary kernels
have protection floors, HOME has floor $p_{ii}$, and every retained NEW
candidate has $p/u$ bounded below by the broad-candidate gate.
Allocate total estimator and coupling failure probability at most
\[
 p_0=\xi_0/[10^6(M_f+\Lambda_f+1)].
\]
Conditional forest kernels preserve entire unconditional row laws, so
base-row error is charged before conditioning, with no division by a rare
parent-bit probability. Direct dyadic categorical kernels introduce no
forbidden transition; make their total variation error at most $p_0$.
All estimator blocks and transitions have a deterministic polynomial bound.
Thus failure and coupling loss is at most
$10p_0(M_f+\Lambda_f)$ per agent. The deterministic losses, including
$64\sqrt\eta$ and even a separately charged reward error
$128C_fd\sqrt\eta$, total at most $101\epsilon_0$ per agent.
Their sum is strictly below $\xi_0$.

All LPs and flows are rational with polynomial bit length. In forest
rounding, derivative-ratio denominators are bounded by downward resets;
in a row program, charge bit-length additions to the elimination tree of
disappearing coordinates. The conditional kernels use direct probabilities,
so they need no rejection from rare events. Maximum-product matchings of
rational values are exact: if edge values have at most $B$ bits, distinct
products of $n$ edges have log separation greater than $2^{-4nB}$.
Log approximations within $2^{-4nB}/(4n)$ therefore suffice for an exact
choice. These observations give worst-case polynomial work, oracle calls
and random bits under the displayed precision schedule.

\begin{theorem}[Finite expected-log guarantee]\label{simp:finite-theorem}
For every rational $\xi>0$, the specified algorithm satisfies
\[
 \mathbb E\log\NSW(Z)\ge\log\OPT-U_\star-\xi,
 \qquad \mathbb E\NSW(Z)\ge e^{-U_\star-\xi}\OPT.
\]
\end{theorem}
\begin{proof}
All scalar cases have been proved, the one-flow repair realizes its
ledger on one common outcome, and \eqref{simp:iterate-bound} plus radial
regret gives the exact inequality. The preceding budget bounds every
finite error by $n\xi$. Jensen's inequality proves the second assertion.
\end{proof}
For $0<\varepsilon\le1$, take $\xi\le\varepsilon/12$.
Since $R_\star<2$, this gives expected factor at most $R_\star+\varepsilon$.
For $Y=\log(\OPT/\NSW)\ge0$ and $a=\log(R_\star+\varepsilon)$,
$a-\mathbb EY\ge\varepsilon/4$ and $a<3/2$; hence Markov's inequality
implies $\Pr(Y<a)\ge\varepsilon/6$.
Repeat $\lceil(6/\varepsilon)\log(1/\delta)\rceil$ times independently
and retain the largest actual utility product. This proves the probability
claim in Theorem~\ref{simp:main-theorem}. For larger $\varepsilon$, use
the algorithm with tolerance $\varepsilon=1$.
The following self-contained lower bound completes its infimum assertion.
